% FICHIER GENERE PAR construire_v3.py, NE PAS EDITER A LA MAIN.
% Source : chapitres/12_resultats.tex, moins les sections
% retirees, qui restent integralement dans main_v2. Toute correction
% se fait dans le fichier source puis par relance du script, sinon
% les deux divergent.
\chapter{Résultats}
\label{chap:resultats}

Ce chapitre livre le \og combien \fg{} de la question du mémoire, sous la discipline posée
aux chapitres~\ref{ch:identifiabilite} et~\ref{chap:partielle} : les niveaux sont
illustratifs, les écarts et les ordres sont ce qui résiste. Tous les résultats sont ancrés,
comme dans l'ancien mémoire, sur une entité type de $15\,000$~M\euro{} de provisions, avec
deux périmètres de sévérité : OpRisk (grandes pertes opérationnelles, queue lourde)
et PRC (sévérité plafonnée à $40$~M\euro). Sauf mention contraire, on lit la médiane
et l'intervalle interquartile, la variance de la sévérité étant infinie ($\xi\approx0{,}9$).

\begin{attention}
\textbf{Comment lire les niveaux de ce chapitre.} Les montants qui suivent sont
\emph{illustratifs} et ne doivent pas être lus comme la mesure du capital d'une entité réelle.
Le chapitre~\ref{chap:partielle} le démontre : en ne faisant varier que la loi de sévérité et la
fréquence, c'est-à-dire des \emph{entrées} du modèle qui dépendent d'un périmètre de collecte
non transposable, le niveau se déplace d'un facteur $41$, tandis que les rapports au socle
varient d'au plus $33\,\%$ et que le groupe de tête du classement des piliers ne change pas. Le
tableau~\ref{tab:scretat} en donne d'ailleurs l'illustration immédiate : le même état conforme
vaut $8\,301$~M\euro{} sous OpRisk et $2\,589$~M\euro{} sous PRC, soit un facteur~$3{,}2$ pour
un seul changement de périmètre de sévérité. Ce chapitre établit donc des écarts, des
rapports et des ordres, et c'est sur eux, jamais sur un niveau, que reposent ses conclusions.
\end{attention}

\begin{cle}
\textbf{Réconciliation avec le socle du chapitre~\ref{chap:partielle}.} Ce chapitre parle d'un
SCR \emph{conforme} d'environ $5\,900$~M\euro{}, quand le chapitre~\ref{chap:partielle} parle
d'un \emph{socle} de $5\,275$~M\euro{}. Ce ne sont pas deux mesures de la même chose, mais deux
points d'un seul modèle : le socle est pris à contagion nulle ($W=0$), là où l'état
conforme en conserve une. L'ordre socle $<$ conforme $<$ non conforme est donc cohérent, et le
pont terme à terme comme l'écart résiduel de Monte-Carlo sont établis au
\S\ref{sec:ann-reconciliation}.
\end{cle}

% SOURCES-SCRIPTS: 65 66 12 16 63

\section{SCR par état de conformité}

Le SCR croît de façon monotone de l'état conforme à l'état non conforme
(figure~\ref{fig:multietats}). Le tableau~\ref{tab:scretat} donne les trois lectures emboîtées :
la lecture~A (fréquence seule) est directement comparable au mémoire ; la lecture~B ajoute la
propagation de la cascade ; la lecture~C ajoute le canal détection.

\begin{table}[htbp]\centering\small
\renewcommand{\arraystretch}{1.25}
\begin{tabular}{@{}l ccc c ccc@{}}
\toprule
& \multicolumn{3}{c}{\textbf{OpRisk} (M\euro)} & & \multicolumn{3}{c}{\textbf{PRC} (M\euro)} \\
\cmidrule(lr){2-4}\cmidrule(lr){6-8}
État & A & B & C & & A & B & C \\
\midrule
Conforme               & 8301  & 6664  & 5932  & & 2589 & 1903 & 1642 \\
Partiellement conforme & 11059 & 9625  & 9625  & & 3653 & 3118 & 3118 \\
Non conforme           & 15074 & 15074 & 16595 & & 5520 & 5520 & 6577 \\
\bottomrule
\end{tabular}
\caption{SCR par état, trois lectures (A~: fréquence ; B~: $+$~propagation ; C~:
$+$~détection). Le canal détection ne joue qu'aux extrêmes (il
% SOURCES-SCRIPTS: 16
abaisse l'état conforme et relève le non conforme), l'état partiel n'étant pas modulé.}
\label{tab:scretat}
\end{table}

\subsection*{Le chiffre de tête : quatre lectures emboîtées et la lecture publiée}
\label{sec:chiffre-de-tete}

Le tableau qui précède n'est pas le seul protocole du mémoire à chiffrer \og l'écart entre l'état
conforme et l'état non conforme \fg{}. Quatre le font, et ils donnent quatre résultats. Aucun
n'est faux : ils ne relâchent pas le même nombre de canaux entre les deux états, et rien
d'autre ne les sépare. Les publier sans le dire reviendrait à laisser le lecteur trancher sans les éléments.

\begin{center}\footnotesize
\renewcommand{\arraystretch}{1.25}
\begin{tabular}{@{}l c r r r@{}}
\toprule
\textbf{Lecture} & \textbf{canaux} & \textbf{conforme} & \textbf{non conf.} & \textbf{facteur} \\
\midrule
Score de conformité seul, $g$ figé au niveau NC & 1 & 7987 & 16\,847 & $2{,}11$ \\
Fréquence $+$ propagation (lecture~B) & 2 & 6664 & 15\,074 & $2{,}26$ \\
$+$ détection (lecture~C) & 3 & 5932 & 16\,595 & $2{,}80$ \\
\textbf{$+$ accumulation P4 (les quatre canaux)} & \textbf{4} & \textbf{6049} & \textbf{20\,188}
& $\mathbf{3{,}34}$ \\
\bottomrule
\end{tabular}
\end{center}

\vspace{4pt}
\noindent En M\euro, périmètre OpRisk. Les niveaux sont ceux des sorties versionnées ; les
\emph{rapports}, eux, sont recalculés à chaque exécution plutôt que
calculés à la lecture, faute de quoi un niveau amont pourrait bouger sans que le facteur suive.
% SOURCES-SCRIPTS: 25 16 43 67

\begin{cle}
\textbf{Le mémoire publie la lecture à quatre canaux : $6\,049 \to 20\,188$~M\euro, soit un
facteur $3{,}34$ et un écart de $14\,139$~M\euro.} Le motif n'est pas qu'elle donne le chiffre
le plus grand, c'est qu'elle est la seule où l'état conforme est conforme sur tous les
canaux que le modèle possède. Dans la première lecture, l'entité dite conforme propage encore
comme une entité défaillante ; dans les deux suivantes, sa détection ou son accumulation tiers
restent au niveau non conforme. Une entité conforme sur un canal et non conforme sur trois autres
n'est pas un état conforme, c'est une \emph{remédiation partielle}, et la nommer conforme
sous-estime l'écart par construction.

\textbf{Ce qui a changé depuis, et qui permet de trancher.} La version précédente de cet encadré
excluait le canal détection tant que son ampleur n'était pas validée. Elle l'est : il pèse
$2\,928 \pm 343$~M\euro{} en valeur de Shapley et
$4\,562 \pm 819$ en marginal de fermeture, signe résolu sur seize graines.
% SOURCES-SCRIPTS: 68 La réserve qui
justifiait l'attente a donc expiré.

\textbf{Les trois autres lectures ne sont pas retirées}, et c'est délibéré : ce sont des
remédiations \emph{partielles}, ce qui est un objet utile en soi, et leur emboîtement chiffre ce
que coûte de ne fermer qu'une partie des canaux.
\end{cle}

\begin{attention}
\textbf{Le facteur se compare d'une lecture à l'autre, l'écart en euros non}, et il faut le dire
plutôt que de laisser le lecteur le remarquer seul : l'écart passe de $8\,860$ à $8\,410$~M\euro{}
entre les deux premières lignes alors que le facteur monte. La cause n'est pas un effet de modèle
mais un changement de base, l'état conforme passant de $7\,987$ à $6\,664$~M\euro{} d'un
protocole à l'autre. Un écart en euros ne se compare qu'à base égale ; un rapport est sans
dimension. C'est le même motif qui fait publier des élasticités plutôt que des plages au
\S\ref{sec:sensibilite-propagation}.

\textbf{Et la monotonie du facteur s'observe, elle ne se teste pas ici.} Les quatre lignes ne
forment pas une expérience contrôlée, leurs protocoles différant aussi par la graine et par la
base. Elle est \emph{cohérente} avec la super-additivité mesurée sur une autre partition, où
relâcher un canal de plus ajoute son effet propre et les croisés qu'il porte, mais le test de
cette super-additivité est ailleurs, dans les seize configurations des quatre canaux, qui
partagent le même moteur et les mêmes graines.
% SOURCES-SCRIPTS: 68
\end{attention}

\noindent\textbf{Où chaque lecture continue de servir.} Les trajectoires du
\S\ref{sec:res-traj} et la priorisation du \S\ref{sec:res-prio} restent conduites en
\emph{lecture~B}. Ce n'est pas une indécision résiduelle mais un choix d'objet : elles portent
sur l'\emph{ordre} de remédiation, dont le chapitre~\ref{ch:robustesse} montre qu'il est
invariant aux leviers testés, de sorte que le canal ajouté déplacerait leurs niveaux sans
déplacer leur conclusion. Les rejouer relèverait par ailleurs de la recalibration, que le gel
interdit.

\begin{figure}[htbp]\centering
\figover{S10_scr_multi_etats.png}
\caption{SCR par état de conformité (trois canaux emboîtés) et bascule des états en crise
systémique.}
% SOURCES-SCRIPTS: 16
\label{fig:multietats}
\end{figure}

Le surcoût de non-conformité $\Delta_{\text{DORA}}=\mathrm{SCR}(\text{NC})-\mathrm{SCR}(\text{C})$,
à graine Monte-Carlo commune (l'écart ne vient que du changement d'état, non du bruit),
s'obtient par bootstrap deux niveaux (multiplicateurs $+$ sévérité). Il ressort du même ordre
que celui du mémoire, amplifié par la propagation, et $100\,\%$ des tirages sont
positifs.

\begin{table}[htbp]\centering\small
\renewcommand{\arraystretch}{1.25}
\begin{tabular}{@{}l cc c c@{}}
\toprule
& médiane & IC90\,\% & & mémoire (réf.) \\
\midrule
$\Delta_{\text{DORA}}$ NC vs C (OpRisk) & 7401 & $[2830\,;38793]$ & & 3879 \\
$\Delta_{\text{DORA}}$ PC vs C (OpRisk) & 2556 & $[963\,;12768]$  & & --- \\
$\Delta_{\text{DORA}}$ NC vs C (PRC)    & 3668 & $[3201\,;4094]$  & & 2015 \\
$\Delta_{\text{DORA}}$ PC vs C (PRC)    & 1257 & $[1077\,;1466]$  & & --- \\
\bottomrule
\end{tabular}
\caption{Surcoût de non-conformité en euros (M\euro), lecture B (fréquence $+$ propagation),
bootstrap deux niveaux. Côté OpRisk la sévérité de queue est ré-échantillonnée sur les
$91$ excès réels, selon la méthode du mémoire, non tirée dans une approximation
% SOURCES-SCRIPTS: 14
normale~; l'intervalle reste large parce que la queue l'est, non par artefact de méthode.
Le surcoût intermédiaire (PC vs C) est chiffré pour la première fois.}
% SOURCES-SCRIPTS: 16
\label{tab:delta}
\end{table}

\noindent \textbf{Lecture actuarielle.} L'intervalle très large côté OpRisk n'est pas un défaut
du modèle : c'est la signature d'une queue d'indice $0{,}60$, où le niveau du capital est
intrinsèquement imprécis. L'écart typique (interquartile) est bien plus resserré que
l'intervalle à $90\,\%$, et le \emph{signe} du surcoût, lui, est certain. On retrouve la thèse du
mémoire : le SCR est une distribution, pas un point.

\paragraph{L'intervalle PRC est étroit, et il le reste quand on cesse de croire la conversion
exacte.} La sévérité PRC est dérivée du nombre d'enregistrements compromis par la relation
log-log de Jacobs (2014), dont le résidu vaut $0{,}523$ en logarithme naturel : à volume donné,
le coût réel varie d'un facteur $2{,}4$ autour de la droite à $90\,\%$. Réinjecter ce résidu
% SOURCES-SCRIPTS: 21
enregistrement par enregistrement ne rouvre pourtant presque pas l'intervalle, le bruit se
moyennant dans l'ajustement GPD ; il déplace le \emph{niveau} du surcoût de $7\,\%$, sans
l'élargir (\S\ref{sec:ann-jacobs}). L'étroitesse de l'intervalle PRC n'est donc pas un artefact
de la conversion déterministe.

\section{Décomposition par pilier}

En basculant un pilier à la fois de conforme à non conforme, le surcoût se décompose
(figure~\ref{fig:decomp}). Le classement obtenu, $P1>P4>P2>P3>P5$, reproduit exactement
le classement ROOT du modèle qualitatif, établi indépendamment par jugement d'expert : deux
voies de construction distinctes convergent sur la même hiérarchie de piliers.

\begin{figure}[htbp]\centering
\figover{S11_decomp_par_pilier.png}
\caption{Décomposition du $\Delta_{\text{DORA}}$ par pilier et validation croisée avec ROOT.
}
% SOURCES-SCRIPTS: 16b
\label{fig:decomp}
\end{figure}

\begin{cle}
\textbf{Une interaction super-additive, établie là où l'estimateur est fiable.} Les surcoûts
isolés ne somment pas au surcoût total : il reste un terme d'\emph{interaction}, signature de la
cascade (un pilier non conforme amplifie aussi les contagions amorcées ailleurs qui le
traversent), qu'un modèle sans propagation (le LDA à copule du mémoire) ne peut pas produire.
Sa mesure exige toutefois de la prudence : une interaction est une différence de différences,
l'estimateur Monte-Carlo le plus fragile qui soit en queue lourde. Sur la VaR à $99{,}5\,\%$
OpRisk, son \emph{signe n'est pas résolu} : sur six graines à résolution égale
($60\,000$ années chacune), il va de $-137$ à $+1\,395$~M\euro{} et n'est positif que dans
quatre cas sur six, la graine de référence donnant $+1\,395$~M\euro{}. Un estimateur
qui change de signe avec la graine interdit d'en tirer quoi que ce soit. La super-additivité
est en revanche robuste partout où le bruit de queue ne domine pas : en \emph{perte moyenne}
elle est positive sur les six graines ($+532$ à $+659$~M\euro{} OpRisk, $+827$
côté PRC, soit environ un cinquième du surcoût moyen), et jusque dans la VaR du périmètre PRC,
dont le plafond de sévérité tronque la queue ($+905$~M\euro{}). Conséquence pratique : puisque les
marginaux ne somment pas au total, l'attribution du surcoût aux piliers exige une règle
principielle qui somme exactement : c'est l'objet de la valeur de Shapley
(section~\ref{sec:allocation}). À ne pas confondre avec l'interaction entre les quatre
\emph{canaux} de la non-conformité, qui décompose le même écart sur une autre partition, vaut
$+5\,001$~M\euro{} et dont le signe, lui, \emph{est} résolu : elle est reprise terme par terme au
\S\ref{sec:kpi}. Les deux ne s'additionnent pas.
% SOURCES-SCRIPTS: 16b 20b 68
\end{cle}

\section{Allocation du capital : Shapley et Euler}
\label{sec:allocation}

La décomposition qui précède pose deux questions d'allocation distinctes. \emph{Qui porte le
surcoût~?} appelle une attribution par pilier \emph{source} de la non-conformité, qui doit
composer avec l'interaction (les marginaux ne somment pas au total). \emph{Où loge le
capital~?} appelle une allocation du niveau $\mathrm{SCR}(\text{NC})$ par pilier
\emph{touché}, cohérente au sens où les parts somment au total. La valeur de Shapley répond
à la première, la règle d'Euler à la seconde (%
% SOURCES-SCRIPTS: 20
figure~\ref{fig:allocation}).

\subsection{Shapley : attribuer le surcoût, interaction comprise}

On munit les piliers d'une fonction caractéristique
\[
v(S) \;=\; \mathrm{SCR}\bigl(S \text{ en NC},\ \text{reste C}\bigr)-\mathrm{SCR}(\text{tous C}),
\qquad S \subseteq \{P1,\dots,P5\},
\]
calculée par la cascade sur les $2^5$ configurations à graine Monte-Carlo commune. Le
marginal de la section précédente est $v(\{k\})$~; la valeur de Shapley
\[
\varphi_j \;=\; \sum_{S \subseteq \{P1,\dots,P5\}\setminus\{j\}}
\frac{|S|!\,\bigl(5-|S|-1\bigr)!}{5!}\,\bigl[v(S\cup\{j\})-v(S)\bigr]
\]
moyenne la contribution de $j$ sur tous les ordres d'arrivée. Par efficience,
$\sum_j \varphi_j = v(\text{tous}) = \Delta_{\text{DORA}}$ \emph{exactement} : la règle
redistribue l'interaction que les marginaux isolés laissent de côté, quel que soit son signe.

La construction est celle de \citet{Shapley1953}, dont l'efficience est la propriété
axiomatique, et son usage actuariel en allocation de capital sous Solvabilité~II est déjà
documenté \citep{Decupere2011}. Les cinq piliers autorisent l'énumération exacte des $5!$
ordres d'arrivée ; les algorithmes par échantillonnage des permutations
\citep{CastroGomezTejada2009, SongNelsonStaum2016}, indispensables en grande dimension, ne
sont donc pas nécessaires ici, et aucune approximation n'entre dans le tableau qui suit.

\begin{table}[htbp]\centering\small
\renewcommand{\arraystretch}{1.25}
\begin{tabular}{@{}l rrr c rrr@{}}
\toprule
& \multicolumn{3}{c}{OpRisk} & & \multicolumn{3}{c}{PRC} \\
\cmidrule{2-4}\cmidrule{6-8}
pilier & marginal & $\varphi_j$ & part & & marginal & $\varphi_j$ & part \\
\midrule
P1 gouvernance & 3481 & 3717 & $34{,}0\,\%$ & & 1375 & 1608 & $32{,}4\,\%$ \\
P4 tiers       & 2658 & 2810 & $25{,}7\,\%$ & & 1075 & 1252 & $25{,}3\,\%$ \\
P2 incidents   & 1634 & 1867 & $17{,}1\,\%$ & &  721 &  920 & $18{,}6\,\%$ \\
P3 tests       & 1094 & 1484 & $13{,}6\,\%$ & &  571 &  738 & $14{,}9\,\%$ \\
P5 partage     &  665 & 1048 &  $9{,}6\,\%$ & &  311 &  438 &  $8{,}8\,\%$ \\
\midrule
somme          & 9531 & 10\,927 & $100{,}0\,\%$ & & 4052 & 4957 & $100{,}0\,\%$ \\
\bottomrule
\end{tabular}
\caption{Valeur de Shapley du surcoût $\Delta_{\text{DORA}}$ (M\euro, NC vs C, trois canaux,
protocole des trois canaux). Les marginaux somment à $9\,531$ (OpRisk), la valeur de
% SOURCES-SCRIPTS: 16b
Shapley à $10\,927$, le surcoût total exactement. Cellules arrondies au million et parts au
dixième de point, échelle à laquelle la colonne somme exactement à cent : c'est la propriété
d'efficacité de la valeur de Shapley, et elle se lit dans le tableau plutôt que de s'y
excuser en note.}
% SOURCES-SCRIPTS: 20
\label{tab:shapley}
\end{table}

Le classement de Shapley reproduit le classement ROOT sur les deux périmètres. Une part de
cette convergence est mécanique : les amorces de la cascade sont semées
proportionnellement à ROOT (parts $30/27/18/15/9\,\%$),
% SOURCES-SCRIPTS: 20
et le canal fréquence d'un pilier
module sa propre part d'amorce. Le test non trivial est la \emph{déformation} : la part de
Shapley de P1 ($34\,\%$) excède sa part d'amorce ($30\,\%$), effet propre de la propagation
(P1 est l'émetteur le plus fort de $W$), et l'interaction redistribuée ($+1\,395$~M\euro{}
OpRisk, $+905$ PRC à cette graine) est intégralement un produit de la cascade, absent d'un
modèle additif.

\subsection{Euler : allouer le niveau, et une limite d'estimation assumée}

La règle d'Euler alloue $\mathrm{SCR}(\text{NC})$ par
$C_j = \mathbb{E}\bigl[L_j \mid L \ge \mathrm{VaR}_{99{,}5\,\%}\bigr]$ (version TVaR, exacte
par construction), où $L_j$ est la perte du pilier \emph{touché} $j$ dans la cascade. Côté
PRC, l'allocation est stable et presque plate : P2 $23\,\%$, P4 $23\,\%$, P1 $20\,\%$,
P3 $18\,\%$, P5 $17\,\%$, et les parts de capital coïncident avec les parts de perte moyenne
parce que le plafond de sévérité tronque la queue
($\mathrm{TVaR} = 6\,795 \approx \mathrm{VaR} = 6\,573$~M\euro). Côté OpRisk, les parts
\emph{moyennes} sont stables (P2 $23\,\%$, P4 $22\,\%$, P1 $20\,\%$, P3 $18\,\%$,
P5 $16\,\%$), mais l'allocation de la queue extrême ne l'est pas, et la cause en est
théorique : avec $\xi \approx 0{,}6 > 0{,}5$, la sévérité est de variance infinie,
l'estimateur Monte-Carlo de la TVaR converge en régime de loi stable, et le classement
conditionnel bascule d'une résolution à l'autre (vérifié à $150$ et $240$~mille années
simulées). On rapporte donc l'allocation de queue du périmètre PRC, régularisée par le cap,
et l'on tient celle d'OpRisk pour indicative.

\paragraph{La mesure de queue sert ici à \emph{mesurer}, jamais à couvrir.} La distinction est
assez souvent négligée pour qu'il faille l'écrire. Le besoin de capital publié dans ce mémoire
est partout un $\VaR_{99{,}5\,\%}(L)$ centré, conformément au niveau de Solvabilité~II ; aucune
grandeur de couverture n'est une TVaR, ni une CTE, ni un Expected Shortfall. La mesure
de queue intervient à deux titres, et tous deux sont des titres de \emph{diagnostic}. Comme
noyau d'allocation d'abord, dans la règle d'Euler ci-dessus, où elle répartit un niveau déjà
fixé sans le modifier. Comme \emph{indicateur de forme de queue} ensuite : le rapport
$\TVaR_{99{,}5\,\%}(L)/\VaR_{99{,}5\,\%}(L) \approx 2{,}1$ se compare à son asymptote théorique
$1/(1-\xi) = 2{,}5$, et c'est ce rapport, non les deux montants, qui porte l'information.
% SOURCES-SCRIPTS: 67
Ce rapport dépend de l'état et il faut
donc dire lequel : $2{,}10$ vaut à l'état \emph{tous non conforme}, celui de cette section, tandis
que l'annexe~\ref{sec:ann-cte} le mesure à $2{,}17$ sur la calibration de référence. Les deux
% SOURCES-SCRIPTS: 73
valeurs sont justes, elles ne portent pas sur la même loi de perte, et l'écart entre elles est
sans signification en soi.

Deux raisons de ne pas en faire un niveau de couverture, et la seconde est décisive. La
première est de proportion : sur un risque aussi large, dont les queues sont plurielles selon la
catégorie d'événement (chapitre~\ref{chap:adaptations}), adosser la couverture à une moyenne de queue
déformerait la distribution et immobiliserait un capital sans commune mesure avec le niveau
réglementaire de référence. La seconde est d'existence : sous la calibration PRC,
$\hat\xi=1{,}033>1$ place la sévérité en régime d'espérance infinie, et la mesure
n'existe alors pas au sens mathématique. Une mesure de couverture qui cesse d'être définie
selon la source de sévérité retenue ne peut pas porter un besoin de capital ; le même objet,
lu comme un diagnostic de queue, reste au contraire parfaitement informatif, puisque sa
divergence \emph{est} l'information.

\begin{cle}
\textbf{Sources et réceptacles.} Les deux allocations ne racontent pas la même chose, et
c'est le résultat : Shapley, vue \emph{source}, concentre le surcoût sur P1 ($34\,\%$), là
où naît la non-conformité~; Euler, vue \emph{réceptacle}, répartit le capital presque
uniformément sur les piliers touchés (P2 et P4 en tête légère). La remédiation se décide
côté source~; le capital, lui, se loge là où les cascades atterrissent.
\end{cle}

\begin{figure}[htbp]\centering
\figover{S15_allocation_shapley_euler.png}
\caption{Allocation du capital DORA (OpRisk). (a)~Shapley contre marginaux : l'interaction
redistribuée. (b)~Euler : parts de capital (queue) contre parts de perte moyenne~; sur ce
périmètre non plafonné, l'allocation fine de la queue extrême reste indicative
($\xi \approx 0{,}6$, variance de sévérité infinie).}
% SOURCES-SCRIPTS: 20
\label{fig:allocation}
\end{figure}

\section{Trajectoire du capital}
\label{sec:res-traj}

Le long de la mise en conformité, le capital décroît de l'état initial vers le SCR conforme, et
le surcoût résiduel $\Delta_{\text{DORA}}(t)$ tend vers zéro (tableau~\ref{tab:traj},
figure~\ref{fig:traj}). Partant de l'état marginal actuel (NC~$35\,\%$, PC~$35\,\%$, C~$30\,\%$),
le capital OpRisk passe de $\sim 10\,750$ à $\sim 7\,040$~M\euro{} sur cinq ans, vers le SCR
conforme ($\sim 6\,925$).

\begin{table}[htbp]\centering\small
\renewcommand{\arraystretch}{1.2}
\begin{tabular}{@{}c ccc@{}}
\toprule
$t$ (ans) & SCR($t$) médiane & bande 90\,\% & $\Delta_{\text{DORA}}(t)$ \\
\midrule
0 & 10752 & $[9092\,;12654]$ & 3827 \\
1 & 9576  & $[7205\,;12040]$ & 2652 \\
2 & 8310  & $[6194\,;11042]$ & 1385 \\
3 & 7575  & $[5924\,;10164]$ & 650 \\
5 & 7041  & $[5857\,;9028]$  & 116 \\
\bottomrule
\end{tabular}
\caption{Trajectoire du capital DORA (M\euro, OpRisk).}
% SOURCES-SCRIPTS: 17
\label{tab:traj}
\end{table}

\begin{figure}[htbp]\centering
\figover{S12_trajectoire_scr.png}
\caption{Trajectoire SCR($t$) et évolution des états de conformité.}
% SOURCES-SCRIPTS: 17
\label{fig:traj}
\end{figure}

\noindent \textbf{Lecture actuarielle.} L'intégrale de la trajectoire, à un coût de capital
près, est le \emph{coût de portage} de la non-conformité pendant la remédiation. La bande reste
large et surtout haute en environnement dégradé : la vraie exposition n'est pas le capital
médian, c'est le risque que la remédiation traîne quand la menace monte, précisément le scénario
où le capital importe le plus.

\section{Priorisation de la remédiation}
\label{sec:res-prio}

Sous budget contraint (un pilier remédié à la fois), l'énumération des $120$ ordres donne
l'ordre optimal $P1\to P4\to P2\to P3\to P5$, égal à l'ordre ROOT et distinct de l'ordre
glouton par valeur d'accélération (figure~\ref{fig:prio}). L'écart de capital intégré entre le
pire ordre et l'optimal atteint $21\,\%$.

\begin{figure}[htbp]\centering
\figover{S13_priorisation.png}
\caption{Valeur d'accélération par pilier et ordre de remédiation optimal.}
% SOURCES-SCRIPTS: 18
\label{fig:prio}
\end{figure}

\begin{cle}
\textbf{Le myope n'est pas l'optimal.} Remédier au coup par coup le pilier au plus fort gain
immédiat (l'ordre glouton) coûte plus cher que l'ordre optimal, parce que les interactions de
cascade imposent de raisonner sur l'horizon entier. Que l'optimum coïncide avec le classement
ROOT, obtenu par une voie qualitative indépendante, est la troisième convergence
euro~$\leftrightarrow$~ordinal du mémoire, après la stabilité de la progéniture
(\S\ref{sec:branch}) et la décomposition par pilier.
\end{cle}

\paragraph{La question réciproque donne un autre classement, et c'est celui qu'un plan de
remédiation utilise.} Tout ce qui précède fixe un état et lit un capital. L'exercice de
résistance \emph{inverse} fixe le capital et cherche les états qui le produisent, et c'est sous
cette forme qu'un dispositif ORSA emploie un modèle. La monotonie établie au \S\ref{sec:kpi} le
rend court : l'ensemble des configurations atteignant une cible est croissant, donc entièrement
décrit par ses éléments minimaux, et il n'en compte jamais plus de quatre sur seize. Deux nombres
résument alors tout le treillis. Le canal de fréquence seul atteint $10\,377$~M\euro,
soit $1{,}72$ fois l'état conforme, tandis que les trois autres canaux \emph{réunis} n'atteignent
que $11\,413$, soit $1{,}89$ fois. Toute cible sous le premier seuil se produit donc par un canal
unique, et toute cible au-dessus du second exige la fréquence. Aucun autre canal, seul et
dans sa plage déclarée, n'atteint même $8\,000$~M\euro. La lecture directe hiérarchise les canaux
par leur contribution ; la lecture inverse désigne celui dont la maîtrise \emph{interdit} les
scénarios sévères, et ce n'est pas la même information. Détail, frontières par cible et réserves
au \S\ref{sec:inverse}.
% SOURCES-SCRIPTS: 92

\section{Robustesse}

Le niveau du capital bouge fortement selon les leviers non calibrés (l'indice de queue $\xi$
domine, facteur $\sim 6$), mais la direction ($\Delta_{\text{DORA}}>0$) et la
\emph{priorité} (P1 en tête) tiennent sur tous les réglages testés (figure~\ref{fig:robust}).

\begin{figure}[htbp]\centering
\figover{S14_robustesse_multietats.png}
\caption{Robustesse : le classement tient, le niveau non.}
% SOURCES-SCRIPTS: 19
\label{fig:robust}
\end{figure}

\begin{cle}
\textbf{La recommandation ne dépend d'aucun paramètre non mesurable.} L'ordre de remédiation
optimal ne dépend que des SCR par configuration, donc il est invariant aux taux de
transition de la chaîne de Markov, pourtant les paramètres les moins calibrables. Ces taux
fixent le calendrier de la trajectoire, jamais la séquence recommandée. Le verdict d'ensemble
rejoint celui de toute la démarche : le classement est robuste là où le niveau ne l'est pas.
\end{cle}

\begin{attention}
\textbf{Le périmètre exact de cette robustesse.} Les tests ci-dessus font
varier les paramètres du modèle \emph{à matrice $W$ fixée}, celle du classeur qualitatif. Ils
établissent donc que P1 arrive en tête \emph{si l'on admet cette direction}. Dès lors que l'on
admet ne pas connaître la direction et que l'on balaie l'ensemble admissible
(chapitre~\ref{chap:partielle}), P1 et P4 deviennent pratiquement à égalité : P1 arrive
premier dans $43{,}6\,\%$ des configurations, P4 dans $37{,}7\,\%$, et l'écart de regret maximal entre
% SOURCES-SCRIPTS: 30
les deux n'est que de $1{,}5\,\%$. Les deux énoncés ne se contredisent pas, ils portent sur des
sources d'incertitude différentes. Ce qui est robuste dans les deux cadres, c'est que P1 et P4
devancent nettement P2, P3 et P5, et que P5 n'arrive \emph{jamais} en tête.
\end{attention}

\paragraph{Les hypothèses restantes ont leur axe de stress.} Trois entrées ne figuraient pas au
tornado : les coefficients de la conversion Jacobs, la charge systémique $\gamma$ et l'ancrage
des probabilités d'état. Stressées à $\pm 1{,}645$ écart-type, elles déplacent des niveaux sans
jamais toucher la direction ni le classement : le surcoût PRC reste contenu entre $3\,125$ et
$3\,977$~M\euro{} autour du central, $\gamma$ ne commande que la sévérité de la bascule en
crise, et l'ancrage des probabilités d'état déplace le capital espéré actuel mais laisse
inchangés le SCR par état et le surcoût NC contre C, qui lui sont conditionnels. Le détail du
protocole de stress figure au \S\ref{sec:ann-hypotheses}.
% SOURCES-SCRIPTS: 22 67

\section{La sensibilité au niveau de propagation}
% PONT v3 : la section complete est dans main_v2, version
% longue et non modifiee. Genere par construire_v3.py.
\label{sec:sensibilite-propagation}
La crainte que l'on adresse spontanément à un modèle de cascade est que son résultat
dépende du niveau de propagation retenu. La sensibilité mesurée dit l'inverse, et c'est
un résultat contre-intuitif qu'il faut énoncer tel quel : \textbf{le niveau de propagation
est le moins sensible des leviers}, et la fragilité du chiffre se situe dans l'ajustement
de valeurs extrêmes, pas dans la contagion. L'ignorance sur la direction, elle, ne se
stresse pas mais se gradue, et son coût maximal est connu d'avance. Les tables complètes
sont dans la version longue de ce mémoire.

\section{Comparaison de modèles : la dépendance n'est pas le mécanisme}
\label{sec:benchmark}

Le mémoire a remplacé le LDA à copule par la cascade dirigée. Ce que ce remplacement
apporte se mesure en confrontant la cascade à des jumeaux copule construits pour lui
ressembler (figure~\ref{fig:benchmark}, périmètre OpRisk, graines
% SOURCES-SCRIPTS: 23
Monte-Carlo communes). Cette section compare les modèles sur l'axe de la \emph{structure de
dépendance}, à marges fixées ; la comparaison sur l'axe de la \emph{famille de loi} (la bande
de modèle) est conduite en partie~\ref{part:limites}, section~\ref{sec:trois-etages}, et les
deux se complètent.

\textbf{La photo.} À marges par pilier \emph{identiques} (celles de la cascade, état
conforme) et corrélation calée, la copule gaussienne retombe sur le SCR de la cascade
($+0{,}1\,\%$) : en queue lourde ($\xi \approx 0{,}6$), la queue annuelle obéit au principe
de la perte unique dominante et le \emph{niveau} de base n'identifie pas la dépendance. La
Student ($\nu = 4$) le surestime de $17\,\%$ en forçant des co-extrêmes que le mécanisme ne
produit pas ($2{,}28$ piliers extrêmes par année de queue contre $1{,}12$). Le choix de
copule est donc à la fois non identifiable et conséquent~; la cascade n'a pas ce degré de
liberté : sa dépendance est produite par le mécanisme, pas choisie.

\textbf{Et ce résultat a été répliqué ailleurs, sans lien avec ce travail.} L'objection naturelle
serait qu'une coïncidence à $+0{,}1\,\%$ entre une copule gaussienne et un mécanisme de cascade
tient à un réglage particulier du protocole. Le préprint de cascade climatique cité au
chapitre~\ref{ch:positionnement} \citep{ccrn2026} conduit la même expérience, à marges appariées,
entre une gaussienne, une Student et sa propre cascade : il obtient des primes centrales
pratiquement indiscernables et ne voit la séparation apparaître qu'en \emph{queue}. Même
protocole, même conclusion, sur un autre péril et par une autre équipe. Ce n'est donc pas un
artefact du réglage retenu ici, mais une propriété du principe de la perte unique dominante,
qui rend le niveau central aveugle à la structure de dépendance.

\begin{attention}
\textbf{Le $+0{,}1\,\%$ est une graine, et la coïncidence est plus fragile que le nombre ne le
laisse croire.} La reprise du protocole sur quatre graines et un échantillon assez grand pour
qu'un quantile à $99{,}95\,\%$ ait un sens donne, au même niveau
% SOURCES-SCRIPTS: 80
$99{,}5\,\%$, un écart moyen de $+4{,}8\,\%$ d'étendue $10$ points entre graines. Le
$+0{,}1\,\%$ publié n'est donc pas faux, il est \emph{un tirage} : la conclusion à retenir n'est
pas que la gaussienne retombe exactement sur la cascade, c'est que leur écart ne dépasse
pas son propre bruit. L'énoncé faible est le seul qui tienne, et il suffit à l'argument.

\textbf{Et la séparation en queue ne vaut que pour la Student.} Le même protocole, décliné sur
huit niveaux de quantile de $50$ à $99{,}95\,\%$, marque d'une étoile les seules cases où l'écart
dépasse l'étendue entre graines. Au-delà de $99\,\%$, seule la Student sort de son
bruit, à $+20{,}4$, $+29{,}8$ puis $+31{,}2\,\%$ ; ni la gaussienne ni l'indépendance n'y
parviennent, leurs étendues atteignant $13$ et $31$ points. La raison est structurelle : la
Student est la seule des trois à porter une dépendance de queue \emph{asymptotique}. Le préprint
doit donc être cité sur le protocole, jamais sur l'ordre du résultat en queue, qui
dépend de l'indice de queue de la sévérité et non de l'architecture du modèle.
\end{attention}

\textbf{L'intervention.} Basculer un pilier en non-conformité est une intervention, pas un
conditionnement. Le jumeau LDA-copule exprime les canaux fréquence et détection (les marges
dépendent de l'état de leur pilier), mais trois manques sont structurels.
\emph{(i)}~La propagation n'a nulle part où vivre : le surcoût total ressort à
$6\,760$~M\euro{} contre $10\,927$ pour la cascade, une sous-estimation de $38\,\%$.
% SOURCES-SCRIPTS: 67
\emph{(ii)}~Les marges ne dépendant que de l'état de leur propre pilier et la copule étant
figée, la perte moyenne totale est \emph{exactement} additive : interaction moyenne nulle
par construction ($+0$~M\euro{} mesuré), là où la cascade en produit $+560$~M\euro{} :
l'interaction moyenne mesurée plus haut est une signature falsifiable qu'aucun modèle
copule-marges ne peut reproduire. \emph{(iii)}~La priorité de remédiation impliquée
s'inverse sur P2/P3, et la dépendance ne durcit pas quand la conformité se dégrade, alors
que c'est précisément le canal que DORA fait bouger ($g$ de $0{,}45$ à $0{,}90$).

\begin{cle}
\textbf{Le verdict du benchmark.} La copule peut imiter la photo d'un état~; elle ne
peut imiter ni le surcoût (propagation), ni l'interaction (additivité en moyenne), ni le
durcissement de la dépendance avec la non-conformité. La dépendance n'est pas le
mécanisme : c'en est l'ombre portée.
\end{cle}

\begin{figure}[htbp]\centering
\figover{S18_benchmark_copule.png}
\caption{Benchmark copule (OpRisk). (a)~Intervention : le jumeau copule suit les marges et
sous-estime le surcoût, la cascade suit les sources. (b)~Photo : perte unique dominante en
queue~; la Student force des co-extrêmes que le mécanisme n'a pas.}
% SOURCES-SCRIPTS: 23
\label{fig:benchmark}
\end{figure}

\section{Mise en regard réglementaire}

La question finale est celle de la portée : de l'effet de la conformité que le modèle
mesure, quelle fraction les cadres existants savent-ils exprimer ? On confronte donc, sur
le même écart de capital entre profil conforme et non conforme, trois dispositifs.
% SOURCES-SCRIPTS: 27

\begin{table}[htbp]\centering\small
\renewcommand{\arraystretch}{1.25}
\begin{tabular}{@{}l r r@{}}
\toprule
\textbf{Cadre} & \textbf{$\Delta$ capital C$\to$NC} & \textbf{part de l'effet cascade} \\
\midrule
Formule Standard (Solvabilité~II) & $0$~M\euro{} & $0\,\%$ \\
LDA-copule (dépendance symétrique) & $6\,760$~M\euro{} & $62\,\%$ \\
Cascade dirigée & $10\,927$~M\euro{} & $100\,\%$ (référence) \\
\bottomrule
\end{tabular}
\caption{Fraction de l'effet DORA que chaque cadre sait exprimer.}
% SOURCES-SCRIPTS: 27 67
\label{tab:benchmark-sf}
\end{table}

\textbf{La Formule Standard est structurellement aveugle.} La charge de risque opérationnel
$\mathrm{SCR}_{\text{op}} = \min(0{,}3\,\mathrm{BSCR}\,;\,0{,}03\,\max(\text{primes},
\text{provisions}))$ ne dépend que du volume : elle est \emph{identique} pour une entité
exemplaire et pour une entité défaillante sur DORA. Sa réponse à la conformité est nulle par
construction. La copule en capte $62\,\%$ (les canaux fréquence et détection) mais rate la
propagation, faute de direction. Seule la cascade exprime l'effet complet. La valeur du
modèle se lit donc en effet DORA rendu visible, non en un niveau de SCR à
interpréter dans l'absolu.

\begin{cle}
\textbf{Jusqu'où l'open data descend, et où il s'arrête.} Un test sur
% SOURCES-SCRIPTS: 28
la base VCDB montre que les sous-piliers de P2 (familles d'action) et, plus grossièrement,
la détection de P3 sont exploitables sur donnée ouverte (plusieurs centaines d'incidents en
finance, $45\,\%$ de détection lente). Mais le pilier des tiers P4 s'effondre : $56$
incidents à acteur partenaire en finance, et la donnée ouverte sous-code notoirement
l'externalisation. Le pilotage fin de P4 exige donc une donnée que seul un registre
d'externalisation, tenu par l'entité ou le régulateur, peut fournir. C'est le même verdict
que celui du chapitre~\ref{ch:identifiabilite}, décliné au niveau des sous-piliers.
\end{cle}

\begin{figure}[htbp]\centering
\figover{W_benchmark_sf.png}
\caption{Ce que chaque cadre exprime de l'effet DORA : $0\,\%$ pour la Formule Standard,
$62\,\%$ pour la copule, $100\,\%$ pour la cascade.}
% SOURCES-SCRIPTS: 27
\label{fig:benchmark-sf}
\end{figure}

\section{Du secteur à l'entité : ce que la descente d'échelle permet}
% PONT v3 : la section complete est dans main_v2, version
% longue et non modifiee. Genere par construire_v3.py.
\label{sec:ancrage}
Les niveaux établis jusqu'ici valent à l'échelle du secteur. Les transposer à une entité
demande une descente d'échelle, dont les deux canaux, la fréquence et la sévérité, sont
estimés séparément et dont aucun n'a une élasticité de un à la taille. Cette descente a
une \textbf{borne inférieure de validité}, publiée, en dessous de laquelle le modèle
affirmerait qu'une entité de quelques milliards subit à peu près la sévérité d'une
institution mondiale. Le protocole, les élasticités et la borne sont dans la version longue de ce mémoire.

\begin{table}[htbp]\centering
\small
\begin{tabular}{lrrrr}
\toprule
Lecture & $\lambda$ (an$^{-1}$) & sévérité & SCR $99{,}5\,\%$ & Perte moyenne \\
\midrule
Secteur financier mondial & $21{,}56$ & $\times 1$ & $8\,123$~M\euro & $983$~M\euro \\
Seuil \og{}$\geq 10$ événements\fg{} & $0{,}210$ & $\times 1$ & $488$~M\euro & $9{,}7$~M\euro \\
Seuil \og{}toutes firmes\fg{} & $0{,}100$ & $\times 1$ & $229$~M\euro & $4{,}6$~M\euro \\
$\lambda$ lu à la taille de la cible & $0{,}092$ & $\times 1$ & $198$~M\euro & $4{,}3$~M\euro \\
\textbf{Les deux canaux composés} & $\mathbf{0{,}092}$ & $\mathbf{\times 0{,}854}$ &
$\mathbf{169}$~M\euro & $\mathbf{3{,}6}$~M\euro \\
\bottomrule
\end{tabular}
\caption{Le même modèle lu à plusieurs échelles. La matrice $W$ et le moteur de cascade sont
identiques partout ; seules la fréquence d'amorce et l'échelle de sévérité changent. Les deux
dernières lignes sont la lecture cohérente, où la fréquence \emph{et} la sévérité sont prises
à la taille de l'entité visée.}
% SOURCES-SCRIPTS: 60
\label{tab:echelle}
\end{table}

\section{Quatre bilans réels, et le statut du chiffre obtenu}
% PONT v3 : la section complete est dans main_v2, version
% longue et non modifiee. Genere par construire_v3.py.
\label{sec:entites-reelles}
La méthode est appliquée à quatre bilans d'entités réelles, anonymisées en classes de
taille, à partir de leurs seuls états publiés. Le point à retenir n'est pas le montant
obtenu mais son statut : il s'agit d'une \textbf{borne supérieure d'ordre de grandeur et non
d'une mesure}, l'entité notionnelle tombant sous la borne inférieure de validité rappelée
ci-dessus. Le détail des quatre bilans, la table de correspondance des champs et le
motif de la requalification sont dans la version longue de ce mémoire.

\section{Ce que l'agrégation réglementaire fait de cet écart}
\label{sec:agregation-scr}

Les sections qui précèdent chiffrent un besoin et un écart. Il reste à dire ce que le capital
total de l'entité en retient, et la réponse ne va pas de soi. Sous Solvabilité~II, le risque
opérationnel n'entre pas dans la matrice de corrélation du \textsc{bscr} : il s'ajoute,
selon l'identité $\mathrm{SCR} = \mathrm{BSCR} + \mathrm{Adj} + \mathrm{SCR}_{Op}$ posée au
chapitre~\ref{chap:reglementaire}. Un besoin de nature opérationnelle porté en \textsc{orsa} se
transmet donc au capital total sans aucun bénéfice de diversification, et cela vaut du
niveau comme de l'écart entre états. Un euro d'écart \textsc{dora} coûte un euro de capital.

\paragraph{Le contrefactuel, et ce qu'il mesure.} Pour dire ce que cette convention implique, on
agrège le même montant $D$ comme un module de plus, corrélé à $\rho$ avec le \textsc{bscr} pris
comme un bloc. Le taux de reconnaissance vaut alors
\[
  \tau(\rho, d) \;=\; \frac{\sqrt{1 + 2\rho d + d^{2}} - 1}{d},
  \qquad d = \frac{D}{\mathrm{BSCR}},
\]
et son développement en $d = 0$ donne $\tau(\rho, d) = \rho + \tfrac{d}{2}\,(1 - \rho^{2}) +
O(d^{2})$. Un montant petit devant le \textsc{bscr} est donc reconnu à hauteur de $\rho$,
et de $\rho$ seulement, quand la convention additive le reconnaît en entier. L'écart entre les
deux traitements vaut $1 - \rho$ : il ne dépend pas de la taille, et il ne se referme qu'à
corrélation parfaite. Ce contrefactuel n'est pas ce que fait Solvabilité~II ; il sert à
\emph{mesurer} la convention en vigueur, jamais à en proposer une autre.

\paragraph{Ce que la donnée publiée suffit à établir.} Le \textsc{bscr} ne figure pas parmi les
champs retenus, et il n'a pas besoin d'y figurer : $\mathrm{Adj} \le 0$ donne
$\mathrm{BSCR} \ge \mathrm{SCR} - \mathrm{SCR}_{Op}$, et le plafond donne
$\mathrm{BSCR} \ge \mathrm{SCR}/1{,}3$. Minorer le \textsc{bscr} \emph{majore} $\tau$ : les taux
ci-dessous sont donc des majorants, et l'écart entre les deux conventions vaut au moins ce qu'on
y lit.

\begin{center}\small
\renewcommand{\arraystretch}{1.25}
\begin{tabular}{@{}l r r r r@{}}
\toprule
\textbf{Entité} & \textbf{Écart C$\to$NC} & \textbf{$d = D/\mathrm{BSCR}$} &
\textbf{$\tau$ à $\rho = 0{,}25$} & \textbf{additif / corrélé} \\
 & M\euro & & & \\
\midrule
Assureur non-vie A & $34{,}9$ & $0{,}0954$ & $29{,}4\,\%$ & $3{,}41$ \\
Assureur non-vie B & $44{,}4$ & $0{,}0575$ & $27{,}7\,\%$ & $3{,}62$ \\
Assureur vie C & $62{,}7$ & $0{,}0514$ & $27{,}4\,\%$ & $3{,}65$ \\
Assureur vie D & $86{,}0$ & $0{,}0076$ & $25{,}4\,\%$ & $3{,}94$ \\
\bottomrule
\end{tabular}
\end{center}

\vspace{2pt}
\noindent À $\rho = 0{,}25$, la convention additive reconnaît ainsi de $3{,}4$ à $3{,}9$ fois ce
qu'un traitement corrélé retiendrait, et le rapport reste dans cette bande étroite sur un facteur
$134$ de taille de bilan, précisément parce que $\tau$ tend vers $\rho$. Lu sur le niveau plutôt
que sur l'écart, où $d$ est plus grand, le facteur va de $2{,}6$ à $3{,}8$. Rapporté au
\textsc{scr} publié de chaque entité, l'écart entre états pèse de $0{,}58\,\%$ pour la plus grande
à $8{,}21\,\%$ pour la plus petite, et ces points de \textsc{scr} sont acquis en entier.

\paragraph{Le plafond, qui ajoute une borne à la cécité déjà établie.} La charge forfaitaire est
plafonnée à $0{,}3 \times \mathrm{BSCR}$, quelle que soit l'exposition opérationnelle réelle. Sur
l'assureur non-vie~A, dont l'état \textsc{s.25.01} publie un module opérationnel de
$59{,}0$~M\euro, ce plafond vaut $109{,}8$~M\euro{} et la marge restante $50{,}8$~M\euro. Le
besoin calculé, $112{,}1$~M\euro, en représente $2{,}21$ fois, et l'écart entre états,
$34{,}9$~M\euro, en représente $0{,}69$ fois. Autrement dit, même un forfait rendu
sensible à la conformité saturerait avant d'avoir exprimé ce besoin : la Formule Standard n'est
pas seulement insensible, elle est bornée. La lecture $\mathrm{Adj} = 0$ est ici \emph{posée et
déclarée}, un ajustement strictement négatif relevant le \textsc{bscr}, donc le plafond, donc la
marge ; l'état \textsc{s.25.01} de l'entité permet de la lever au cas par cas.

\begin{cle}
\textbf{La lecture décisionnelle, et elle renverse une intuition de gestion.} Un euro économisé
sur la non-conformité \textsc{dora} vaut, en capital, davantage qu'un euro économisé sur un risque
qui se diversifie : le premier est reconnu en entier, le second à hauteur de sa corrélation. Une
direction des risques qui classerait ses chantiers de remédiation au seul vu de la sinistralité
attendue sous-estimerait donc le rendement en capital de celui-ci, d'un facteur de l'ordre de
trois à quatre. Cette conclusion ne dépend d'aucun paramètre du modèle de cascade : elle tient à
la place que le règlement assigne au risque opérationnel, et elle survivrait à un changement
complet de calibration.
\end{cle}

% SOURCES-SCRIPTS: 95 65

\section{Détenir ou transférer, en deux phrases}
% PONT v3 : la section complete est dans main_v2, version
% longue et non modifiee. Genere par construire_v3.py.
\label{sec:detenir-transferer}
Le capital libéré par une remédiation peut être comparé au prix de marché du même risque.
Deux choses en ressortent, et une seule est une sortie du modèle. La valeur présente de
l'économie de portage vaut exactement la variation de capital, quel que soit le taux
retenu : la révision du coût du capital déplace le nombre d'années, pas l'économie. Et
le bénéfice a deux composantes dont une seule est monétisée, la sinistralité évitée
pesant plusieurs fois le portage, si bien que le retour calculé sur le seul portage est
un \textbf{majorant} et non un plancher. Le calcul complet est dans la version longue de ce mémoire.

